\abstract % Структурный элемент: РЕФЕРАТ

\keywords{ДЕТСКИЕ РИСУНКИ, ПСИХОЛОГИЧЕСКАЯ ДИАГНОСТИКА, КЛАССИФИКАЦИЯ ИЗОБРАЖЕНИЙ, СВЁРТОЧНЫЕ НЕЙРОННЫЕ СЕТИ, MULTI-TASK LEARNING, ОБРАБОТКА ИЗОБРАЖЕНИЙ, КОМПЬЮТЕРНОЕ ЗРЕНИЕ, ТЕСТ ВЕНГЕРА, ПСИХОЛОГИЧЕСКИЙ ПОРТРЕТ}

Объектом разработки в данной работе является программный комплекс автоматического анализа детских рисунков человека для оценки психологических характеристик ребёнка.

Цель работы --- разработать нейросетевую систему распознавания визуальных признаков детских рисунков и их сопоставления с психологическими чертами на основе методики А.~Л.~Венгера.

Основное содержание работы состояло в проектировании собственной двухпотоковой свёрточной нейронной сети, разработке алгоритма предобработки изображений на основе классических операций обработки изображений, обучении модели на 117 признаках в режиме multi-task и реализации модуля генерации текстового психологического портрета. Для удобства практического использования разработанные модули объединены в микросервисный веб-сервис с REST API и веб-интерфейсом.

Основным результатом работы является программный комплекс на языке Python с использованием библиотек PyTorch и OpenCV, оформленный в виде набора Docker-контейнеров.

Данные результаты разработки предназначены для использования специалистами в области детской психологии в качестве вспомогательного инструмента предварительного анализа, а также в исследовательских целях для анализа больших объёмов рисуночных тестов.

Внедрение разработки позволяет автоматизировать рутинную часть обработки рисуночных тестов, что снижает нагрузку на специалиста и повышает воспроизводимость результатов.
